\section{Related work and the nearest-prior boundary}
\label{sec:related}

\paragraph{Intervention value at checkpoints.}
\citet{calibrationnotcontrol2026} names intervention advantage as the decision
quantity, formalises the prediction/control mismatch as target error, and
introduces prefix branching: candidate actions are executed from an identical
trajectory state, with reconstruction checked by step-by-step agreement and
mismatches discarded, and a repeated shallow-prefix probe on an exact-replayable
subset.  The inspected version records a fixed action menu and no pre-treatment
availability record; its cost accounting and harm endpoints were not read in the
available scope.  This paper does not contest that prefix branching measures
intervention value.  It asks the allocation question that remains after
feasibility, cost and harm are made first-class and observed before treatment,
and it tests fitted policies under matched information rather than the
population ordering.  \citet{interventionparadox2026} show that a critic with
strong offline discrimination can degrade agent performance and decompose the
effect into recovery and disruption; the no-fault stratum here is a designed
endpoint, not a rediscovery of that observation.
\citet{harnessr12026} learn lifecycle-wide executable interventions and report
that many valid patches reduce batch success; we ask when a proposed
intervention should be deployed at a particular checkpoint, under a budget.

\paragraph{Replay and measurement preconditions.}
\citet{causalagentreplay2026} re-execute trajectories under step interventions
to attribute failure, reporting an action-match rate; real side-effecting tools
are out of scope there.  \citet{replaygap2026} fork trajectories by replaying
logged actions into a fresh container with a matched control arm and report
reconstruction rates, with two arms and fixed fork points.  Replay fidelity has
been measured by action-match and reconstruction rates; the measurement
contract this paper assumes --- per-layer restoration, pre-treatment
availability, and an explicit identified set for the intended effect under an
unaudited discrepancy --- is developed separately and cited, not re-derived
here.

\paragraph{Ordinal sufficiency and selective prediction.}
The condition under which an outcome ranking coincides with an effect ranking is
established in budgeted triage \citep{okati2021}, in the uplift literature under
monotone treatment amplification \citep{fernandezloria2022}, and for allocation
at every budget fraction \citep{kamran2024}.  Selective prediction and
risk--coverage rank by confidence to decide whether to abstain; here the ranked
object is a causal contrast between executed options, one of which can harm a
clean checkpoint, and the budget is spent on intervention rather than
abstention.  The exact ordinal result is cited as prior art; what this paper
tests is whether \emph{fitted} policies differ at matched information.

\paragraph{Uplift, dynamic treatment regimes, and safe improvement.}
Uplift modelling and policy learning establish that the ranking needed for
treatment can differ from outcome risk
\citep{rzepakowski2012,kitagawa2018,atheywager2021}; learning to defer supplies
the finite-capacity version \citep{madras2018,mozannar2020}; dynamic treatment
regimes supply the complete-regime semantics
\citep{murphy2003,murphybingham2009}.  The single-checkpoint contrast is a
\emph{blip}, and the one-step restriction keeps the hypotheses checkable rather
than assumed \citep{deepblip2025}.  Pessimistic policy improvement supplies the
held-out safety argument \citep{laroche2019,kumar2020,xie2021}; simultaneous
distribution-free safety over guardrail outcomes is established independently
\citep{snpl2025}, and the gate here differs by the availability mask and by the
explicitly narrowed scope of the first study.

\paragraph{What this paper does not claim.}
It does not claim that prediction and intervention value come apart; it does not
claim the ordinal-sufficiency theorem; it does not claim a new general
identification theorem for replay; it does not claim a new safety theorem from a
union bound; and it does not claim real-agent results before the qualification
and pilot gates pass.
